\documentclass[../../../include/open-logic-section]{subfiles}

\begin{document}

\olfileid[tr]{sth}{choice}{intro}

\olsection{Giriş}

\crefrange{sth:cardinals::chap}{sth:card-arithmetic::chap} boyunca
kardinal sayıları sıra sayıları olarak ele alıp bir kardinal sayılar teorisi
geliştirdik. Bu yaklaşım İyi Sıralama Aksiyomu'na dayanır. İyi Sıralama'nın
başka bir ilkeye---Seçim Aksiyomu'na---denk olduğu ortaya çıkar ve bu
ilkenin kabul edilebilirliği üzerine ciddi felsefi tartışmalar yürütülmüştür.
Bu bölümdeki sorularımız şunlardır: Aksiyom nasıl kullanılır ve
gerekçelendirilebilir mi?

\end{document}
